% SPDX-License-Identifier: CC-BY-NC-ND-4.0
% Copyright (c) 2026 Aymix — workbook content. See LICENSE-CONTENT.
% ============================ FRONT MATTER ============================
\newcommand{\frontheading}[2]{%
  \par\addvspace{7mm}%
  \noindent\begin{tikzpicture}[baseline=(t.base)]
    \node[fill=cBrand, text=white, rounded corners=1.3mm, font=\normalsize,
          inner sep=0pt, minimum width=8mm, minimum height=8mm] (n) {#1};
    \node[anchor=west, font=\sffamily\bfseries\LARGE, text=cInk] (t) at (n.east) [xshift=3mm] {#2};
  \end{tikzpicture}\par\vspace{1mm}%
  \noindent{\color{cBrand}\rule{\linewidth}{1pt}}\par\addvspace{2.5mm}%
}

% ---------------------------------------------------------------- COVER
\begin{titlepage}
\thispagestyle{empty}
\begin{tikzpicture}[remember picture, overlay]
  \fill[cBrandDeep] (current page.south west) rectangle (current page.north east);
  % soft geometric accents — indigo title band with amber rules top & bottom
  \fill[cBrand] ([yshift=-4.2cm]current page.north west) rectangle ([yshift=-9.7cm]current page.north east);
  \fill[cAccent] ([yshift=-4.2cm]current page.north west) rectangle ([yshift=-4.05cm]current page.north east);
  \fill[cAccent] ([yshift=-9.85cm]current page.north west) rectangle ([yshift=-9.7cm]current page.north east);
  % faint big glyph, tucked into the lower-right corner
  \node[anchor=south east, text=white, opacity=0.04, font=\bfseries, scale=22]
     at ([xshift=0.6cm,yshift=-0.4cm]current page.south east) {\{\ \}};
  % eyebrow
  \node[anchor=north west, text=cAccent, font=\sffamily\bfseries]
     at ([xshift=2.2cm,yshift=-2.5cm]current page.north west)
     {\Large BACKEND ENGINEERING\, \textbullet\, SPEEDRUN SERIES};
  % title block (centered in the indigo band)
  \node[anchor=north west, text=white, font=\sffamily\bfseries, align=left]
     at ([xshift=2.2cm,yshift=-5.25cm]current page.north west)
     {{\fontsize{39}{43}\selectfont The Backend Engineer's}\\[3mm]
      {\fontsize{39}{43}\selectfont Lab Notebook}};
  % subtitle — placed BELOW the band, on the dark field
  \node[anchor=north west, text=white, font=\sffamily, align=left]
     at ([xshift=2.2cm,yshift=-10.7cm]current page.north west)
     {\large A guided-practice companion to the};
  \node[anchor=north west, text=cAccent, font=\sffamily\bfseries, align=left]
     at ([xshift=2.2cm,yshift=-11.28cm]current page.north west)
     {\large Java \& Spring Boot --- 14-Day Speedrun Roadmap};
  % tagline
  \node[anchor=north west, text=white!85, font=\sffamily, align=left, text width=15.6cm]
     at ([xshift=2.2cm,yshift=-12.9cm]current page.north west)
     {From \emph{reading} concepts to truly \emph{understanding} them --- then building, debugging,
      securing, testing, containerizing, and shipping a production Spring Boot backend, one guided day at a time.};
  % stat row
  \node[anchor=north west, align=left] at ([xshift=2.2cm,yshift=-15.3cm]current page.north west)
     {\begin{tikzpicture}[x=3.7cm, node distance=0pt]
        \foreach \i/\v/\l in {0/14/DAYS, 1/42/GUIDED LABS, 2/110+/EXERCISES, 3/1/CAPSTONE}{
          \node[anchor=west, text=cAccent, font=\sffamily\bfseries] at (\i,0.35) {\Huge \v};
          \node[anchor=west, text=white!80, font=\sffamily\bfseries\footnotesize] at (\i,-0.35) {\l};
        }
      \end{tikzpicture}};
  % notebook owner lines (workbook touch)
  \node[anchor=north west, text=white!75, font=\sffamily, align=left]
     at ([xshift=2.2cm,yshift=-18.9cm]current page.north west)
     {This notebook belongs to:\ \ \textcolor{white!35}{\rule{7cm}{0.4pt}}};
  \node[anchor=north west, text=white!75, font=\sffamily, align=left]
     at ([xshift=2.2cm,yshift=-19.8cm]current page.north west)
     {Started on:\ \ \textcolor{white!35}{\rule{3.6cm}{0.4pt}}\qquad Target finish:\ \ \textcolor{white!35}{\rule{3.6cm}{0.4pt}}};
  % footer
  \node[anchor=south west, text=white!55, font=\sffamily\footnotesize]
     at ([xshift=2.2cm,yshift=1.4cm]current page.south west)
     {Study the roadmap first \textbullet\ then learn with this workbook open beside you \textbullet\ write in the margins};
\end{tikzpicture}
\end{titlepage}

% ---------------------------------------------------------------- HOW TO USE
\frontheading{\faCompass}{How to Use This Workbook}

This is not an exam book and not a second textbook. It is a \keyterm{lab manual} — a place to
\emph{do the work} beside the roadmap. The roadmap tells you \emph{what} to learn; this workbook
walks you through \emph{understanding} it, \emph{building} it, and \emph{reflecting} on it, so the
knowledge sticks. Each chapter maps one-to-one to a day of the roadmap and always follows the same
eleven-part rhythm, so you always know where you are.

\wbsub{The daily study loop}
\begin{wbfigure}{Fig 0.1 — Run this loop once per day. The learning is in the middle two steps, not the reading.}
\wbscale{\begin{tikzpicture}[node distance=7mm and 12mm, every node/.style={font=\sffamily\footnotesize}]
  \node[wbbox, text width=30mm] (a) {\textbf{1 · Read}\\\scriptsize the roadmap day, first};
  \node[wbboxaccent, right=of a, text width=30mm] (b) {\textbf{2 · Work the 11 parts}\\\scriptsize concepts, diagrams, callouts};
  \node[wbboxgood, right=of b, text width=30mm] (c) {\textbf{3 · Build by hand}\\\scriptsize labs \& ShopFlow, no copy-paste};
  \node[wbboxdb, below=8mm of c, text width=30mm] (d) {\textbf{4 · Reflect \& check}\\\scriptsize write-in spaces, checklist};
  \node[wbboxghost, below=8mm of a, text width=30mm] (e) {\textbf{Next day}\\\scriptsize each chapter builds on the last};
  \draw[wbarrow] (a)--(b); \draw[wbarrow] (b)--(c); \draw[wbarrow] (c)--(d);
  \draw[wbarrow] (d) -| (e); \draw[wbarrow] (e)--(a);
\end{tikzpicture}}
\end{wbfigure}

\wbsub{The eleven parts of every chapter}
\begin{center}\small
\begin{tabularx}{\linewidth}{@{}p{4mm}L{36mm}X@{}}
\textbf{1} & \textbf{The Big Picture} & Why this exists, where it fits, how it connects to earlier days.\\[1.2mm]
\textbf{2} & \textbf{Concept Map} & One picture of the whole territory and its vocabulary.\\[1.2mm]
\textbf{3} & \textbf{Guided Learning} & Each idea taught progressively: what / why / how-internally / when / when-not / production / mistakes.\\[1.2mm]
\textbf{4} & \textbf{Visualizations} & Redraw the key diagrams from memory to lock them in.\\[1.2mm]
\textbf{5} & \textbf{Guided Coding Lab} & Build a feature step by step — the workbook guides, you write the code.\\[1.2mm]
\textbf{6} & \textbf{Thinking Exercises} & Reflection prompts (not tests) that make your reasoning explicit.\\[1.2mm]
\textbf{7} & \textbf{Debugging Practice} & A broken scenario worked as Observe → Hypothesise → Verify → Fix → Prevent.\\[1.2mm]
\textbf{8} & \textbf{Mini-Labs} & Three graded labs: Beginner, Intermediate, Production-style.\\[1.2mm]
\textbf{9} & \textbf{Engineering Notes} & Senior judgment: how experienced engineers actually think.\\[1.2mm]
\textbf{10} & \textbf{Build-Along Project} & Extend \keyterm{ShopFlow}, one production backend built across all 14 days.\\[1.2mm]
\textbf{11} & \textbf{Chapter Summary} & Takeaways, pitfalls, a one-page cheat sheet, and what comes next.\\
\end{tabularx}
\end{center}

\begin{prodtip}
The reflection spaces and "redraw from memory" boxes are meant to be \emph{written in} by hand. Print
the pages you are working, or annotate the PDF. Reproducing a diagram yourself is what moves it from
"I recognise that" to "I can explain that in an interview."
\end{prodtip}

\begin{secwarn}[on copy-paste]
This workbook deliberately does \emph{not} hand you finished solutions. The labs guide you to build them
yourself. If you paste code you did not type and reason through, you will feel productive and learn nothing.
The struggle in Parts 5, 7 and 10 is the point.
\end{secwarn}

% ---------------------------------------------------------------- LEGEND
\frontheading{\faInfoCircle}{Field Guide to the Margin Notes}

Throughout the workbook, coloured callouts flag different \emph{kinds} of thinking. Learn to read them at a glance.

\begin{seniornote}How experienced engineers reason about a decision — the judgment behind the code.\end{seniornote}
\begin{prodtip}A concrete practice that holds up in real production systems.\end{prodtip}
\begin{perfnote}A latency, memory, or throughput consideration worth measuring.\end{perfnote}
\begin{secwarn}A security or correctness trap — read these twice.\end{secwarn}
\begin{dbinsight}Something the database is doing underneath your ORM.\end{dbinsight}
\begin{interview}A question you will likely be asked, and how to answer it well.\end{interview}
\begin{misconception}A belief many beginners hold that is subtly wrong.\end{misconception}
\begin{behind}What the framework or runtime is actually doing behind the scenes.\end{behind}

\smallskip
And the working blocks you will act on:

\begin{minilab}{Beginner / Intermediate / Production}{graded practice}
Self-contained labs with a Goal, Context, Requirements, Hints, Expected outcome, and Common pitfalls.
\end{minilab}
\begin{debuglab}{A broken app to diagnose}
Symptom and evidence, then a guided Observe → Hypothesise → Verify → Fix → Prevent walkthrough.
\end{debuglab}
\begin{thinkbox}
A reflection prompt with ruled space to write your answer by hand.
\end{thinkbox}
\begin{buildalong}[ShopFlow]
The one production backend you extend a little further every single day.
\end{buildalong}

% ---------------------------------------------------------------- THE ARC
\frontheading{\faMapSigns}{The 14-Day Arc at a Glance}

Each day earns its place and feeds the next. Read down the "ShopFlow milestone" column to see one real
backend take shape from an empty project to a production-ready service.

\begin{center}\footnotesize
\wbrowcolors
\begin{tabularx}{\linewidth}{@{}C{7mm}L{34mm}X@{}}
\rowcolor{cBrand}\textcolor{white}{\textbf{Day}} & \textcolor{white}{\textbf{Topic}} & \textcolor{white}{\textbf{ShopFlow milestone you build}}\\
1 & Spring Boot Internals & Base module: layered packages, typed config, profiles, \code{/ping}.\\
2 & Database Mastery & Schema + constraints + composite index + migration.\\
3 & JPA \& Hibernate & Order/OrderItem/Product entities, locking, one-query listing.\\
4 & REST API Design & Order API with DTOs, validation, errors, pagination, OpenAPI.\\
5 & Spring Security & JWT login/refresh, BCrypt, roles, locked-down CORS.\\
6 & Testing & Unit + Testcontainers + MockMvc coverage, incl. auth failures.\\
7 & Clean Architecture & Refactor: validator, domain event, async email listener.\\
8 & Docker & Multi-stage image + Compose (app + Postgres + Redis).\\
9 & Redis \& Caching & Cached catalog reads, evict-on-write, TTL.\\
10 & Messaging (MQ) & OrderPlaced to RabbitMQ, retries, dead letter queue.\\
11 & JVM \& Concurrency & Bounded thread pool, Actuator, custom metrics.\\
12 & CI/CD & GitHub Actions: test, build, tag by SHA, push image.\\
13 & Production Readiness & Circuit breaker, rate limiting, liveness/readiness, Grafana.\\
14 & Capstone & Assemble, verify, and present the whole ShopFlow backend.\\
\end{tabularx}
\end{center}

\begin{keyidea}
\textbf{ShopFlow} is deliberately a \emph{small} domain — roughly two entities of real complexity — carrying
a \emph{large} amount of engineering weight. That is exactly the shape of a strong take-home assignment or a
well-scoped first project at a new job. Depth over breadth, every day.
\end{keyidea}

\clearpage
% ---------------------------------------------------------------- TOC
\frontheading{\faListUl}{Contents}
\vspace{1mm}
{\hypersetup{linkcolor=cInk}%
 \makeatletter\@starttoc{toc}\makeatother}
\vfill
\begin{center}\footnotesize\itshape\color{cInk!60}
Fourteen chapters \textbullet\ one growing project \textbullet\ every concept practised, not just read.
\end{center}
\clearpage
